\section*{Hoofdstuk \ref{ch:g5:water-shapes-the-land} --- Water vormt het land}

\begin{solution}{exo:g5:water-shapes-the-land:1}
(a) \omterm{def:g5:water-shapes-the-land:weathering}{Verwering}: het \omterm{def:g3:stones-pebbles-sand:rock}{gesteente} breekt waar het ligt. (b) \omterm{def:g5:water-shapes-the-land:erosion}{Erosie}: het \omterm{def:g3:stones-pebbles-sand:sand}{zand}
slijt af en wordt weggevoerd. (c) \omterm{def:g5:water-shapes-the-land:deposition}{Afzetting}: het \omterm{def:g3:stones-pebbles-sand:sand}{slib} zakt naar de bodem.
(d) \omterm{def:g5:water-shapes-the-land:weathering}{Verwering}: de wortel breekt de steen waar die staat.
\end{solution}

\begin{solution}{exo:g5:water-shapes-the-land:2}
De \omterm{def:g5:water-shapes-the-land:source-mouth}{bron} is de plek waar de \omterm{def:g1:landscapes-around-us:water}{rivier} begint (een \omterm{def:g5:water-shapes-the-land:source-mouth}{bron} in de grond, een
smeltende gletsjer, een hoog nat gebied). De \omterm{def:g5:water-shapes-the-land:source-mouth}{monding} is de plek waar ze
eindigt, in de \omterm{def:g2:water-on-earth:ocean}{zee}, een \omterm{def:g1:landscapes-around-us:water}{meer} of een grotere \omterm{def:g1:landscapes-around-us:water}{rivier}. Een \omterm{def:g5:water-shapes-the-land:source-mouth}{zijrivier} is een
kleinere \omterm{def:g1:landscapes-around-us:water}{rivier} die uitkomt in een grotere.
\end{solution}

\begin{solution}{exo:g5:water-shapes-the-land:3}
De \omterm{def:g5:water-shapes-the-land:source-mouth}{bron}; het smalle V-vormige \omterm{def:g1:landscapes-around-us:landform}{dal}; de \omterm{def:g5:water-shapes-the-land:meander}{meanders}; de \omterm{def:g5:water-shapes-the-land:delta}{delta}.
\end{solution}

\begin{solution}{exo:g5:water-shapes-the-land:4}
(a) Niet waar: vorst, water, wortels, warmte en kou breken langzaam zelfs
hard \omterm{def:g3:stones-pebbles-sand:rock}{gesteente} (\omterm{def:g5:water-shapes-the-land:weathering}{verwering}). (b) Waar: hoe sneller het water, hoe groter de
stukken die het kan verplaatsen. (c) Niet waar: de zwaarste stukken zakken
het eerst; het fijnste \omterm{def:g3:stones-pebbles-sand:sand}{slib} zakt het laatst, waar het water bijna
stilstaat, in een \omterm{def:g1:landscapes-around-us:water}{meer} of in de \omterm{def:g2:water-on-earth:ocean}{zee}.
\end{solution}

\begin{solution}{exo:g5:water-shapes-the-land:5}
A ligt aan de buitenkant van de bocht: daar is het water snel en diep, en
het slijt de oever weg (\omterm{def:g5:water-shapes-the-land:erosion}{erosie}). B ligt aan de binnenkant van de bocht:
daar is het water traag, en het legt \omterm{def:g3:stones-pebbles-sand:sand}{zand} neer (\omterm{def:g5:water-shapes-the-land:deposition}{afzetting}).
\end{solution}

\begin{solution}{exo:g5:water-shapes-the-land:6}
In de \omterm{def:g1:landscapes-around-us:landform}{bergen} is de \omterm{def:g1:landscapes-around-us:water}{rivier} steil en snel: ze kan grote stenen meerollen, en
ze voert de kleine korrels weg. Dicht bij de \omterm{def:g2:water-on-earth:ocean}{zee} is de \omterm{def:g1:landscapes-around-us:water}{rivier} traag: ze
kan geen grote stenen meer verplaatsen, alleen \omterm{def:g3:stones-pebbles-sand:sand}{zand} en \omterm{def:g3:stones-pebbles-sand:sand}{slib}, en die legt ze
daar neer. Onderweg zijn de stenen bovendien gebroken en afgesleten tot
kleinere korrels.
\end{solution}

\begin{solution}{exo:g5:water-shapes-the-land:7}
(a) Bovenloop: $600 \div 6 = 100$ kilometer. Middenloop:
$600 \div 2 = 300$ kilometer. (b) Benedenloop: $600 - 100 - 300 = 200$
kilometer. Omdat $3 \times 200 = 600$, is dat $\frac{1}{3}$ van de \omterm{def:g1:landscapes-around-us:water}{rivier}.
(Controle: $\frac{1}{6} + \frac{1}{2} + \frac{1}{3} = \frac{1}{6} +
\frac{3}{6} + \frac{2}{6} = \frac{6}{6}$.)
\end{solution}

\begin{solution}{exo:g5:water-shapes-the-land:8}
De lichte plekken zijn \omterm{def:g3:stones-pebbles-sand:sand}{zand} en \omterm{def:g3:stones-pebbles-sand:pebble}{grind} die aan de binnenkant van de bochten
zijn neergelegd, waar het water traag is (\omterm{def:g5:water-shapes-the-land:deposition}{afzetting}). De donkere, gebogen
\omterm{def:g1:landscapes-around-us:water}{meren} zijn oude lussen van de \omterm{def:g1:landscapes-around-us:water}{rivier}: twee lussen groeiden zo dicht naar
elkaar toe dat de \omterm{def:g1:landscapes-around-us:water}{rivier} door de landtong ertussen brak, de korte weg nam
en de oude lus als een \omterm{def:g1:landscapes-around-us:water}{meer} achterliet.
\end{solution}

\begin{solution}{exo:g5:water-shapes-the-land:9}
De vlakke weilanden naast de \omterm{def:g1:landscapes-around-us:water}{rivier} zijn haar \omterm{def:g5:water-shapes-the-land:flood}{overstromingsvlakte}: het
land dat de \omterm{def:g1:landscapes-around-us:water}{rivier} bedekt als ze buiten haar oevers treedt. Ze zijn juist
vlak omdat de \omterm{def:g1:landscapes-around-us:water}{rivier} ze uit haar eigen \omterm{def:g3:stones-pebbles-sand:sand}{slib} heeft opgebouwd. Op een dag
treedt de \omterm{def:g1:landscapes-around-us:water}{rivier} weer buiten haar oevers en loopt de school vol water, en
snel stromend overstromingswater is gevaarlijk. De school moet op hoger
gelegen grond worden gebouwd, en de weilanden moeten akkers, weilanden of
een park blijven.
\end{solution}

\begin{solution}{exo:g5:water-shapes-the-land:10}
(a) Een jaar heeft 12 maanden: $12$ miljard $\div\, 12 = 1$ miljard ton per
maand, gemiddeld. (b) Een miljard is $1\,000\,000\,000$; gedeeld door 10 is
dat $100\,000\,000$: honderd miljoen vrachtwagenladingen.
\end{solution}

\begin{solution}{exo:g5:water-shapes-the-land:11}
Een mogelijk verhaal. De korrel is deel van een rots hoog op een \omterm{def:g1:landscapes-around-us:landform}{berg}.
Nacht na nacht opent de vorst een spleet eromheen, tot er een stuk \omterm{def:g3:stones-pebbles-sand:rock}{gesteente}
afbreekt: dat is \omterm{def:g5:water-shapes-the-land:weathering}{verwering}. Het stuk valt in de beek beneden. In de
bovenloop is het water snel en steil: het rolt het stuk mee (\omterm{def:g5:water-shapes-the-land:erosion}{erosie}), en
door tegen andere stenen te botsen breekt het in \omterm{def:g3:stones-pebbles-sand:pebble}{kiezels} en daarna in
zandkorrels. In de middenloop gaat de \omterm{def:g1:landscapes-around-us:water}{rivier} langzamer stromen; de korrel
rust een tijdje op het zandstrandje in de binnenbocht van een \omterm{def:g5:water-shapes-the-land:meander}{meander}, tot
een \omterm{def:g5:water-shapes-the-land:flood}{overstroming} hem weer meeneemt. In de benedenloop is de \omterm{def:g1:landscapes-around-us:water}{rivier} traag en
breed; de korrel reist met het water mee naar de \omterm{def:g5:water-shapes-the-land:source-mouth}{monding}. Daar staat het
water bijna stil, en de korrel komt tot rust met het \omterm{def:g3:stones-pebbles-sand:sand}{slib} en het \omterm{def:g3:stones-pebbles-sand:sand}{zand} van
de \omterm{def:g5:water-shapes-the-land:delta}{delta}: dat is \omterm{def:g5:water-shapes-the-land:deposition}{afzetting}.
\end{solution}
